\section{Introduction}

How large can a list of nearby interleaved codewords be, and how often does
scalar agreement extend to a received pair on the same coordinates? We
study these questions for constant-rate Reed--Solomon codes over smooth
evaluation domains, using the list and correlated-agreement formulations of
\cite{arnonBonehFenzi2026}. We work with the two-fold interleaving
\(\Cpair\subseteq(\F^2)^n\).  Every received pair
\(p=(p_0,p_1)\) determines the scalar affine line
\[
  p\in(\F^2)^n
  \quad\longmapsto\quad
  y_\gamma=p_0+\gamma p_1\in\F^n .
\]
List decoding controls the ambiguity of the full pair, while mutual correlated
agreement bounds the fraction of challenges for which scalar agreement fails
to extend to pair agreement on the same coordinate set.

At the capacity-grid radius \((n-k)/n\), we prove the exact universal list
size
\[
  L_m\!\left(\mathcal C,\frac{n-k}{n}\right)=\binom nk
\]
for every positive interleaving width.  A matching next-shell construction
gives the resulting three-regime field-size phase split. For pair MCA, an
injective compressed-certificate map proves the radius-independent bound
\[
  |\Bad_{\MCA}|\le\binom n{k+1}.
\]
An explicit smooth rate-\(1/4\) family over \(\mathbb F_{3^n}\) meets both
\(2^{-128}\) fractional targets. Its actual field/subgroup realization,
endpoint constructions, MCA certificate bound, complete phase split, real
safe intervals and generic distance/calibration endpoints are machine-checked
in Lean. The endpoint arguments are elementary interpolation; the exponential
alphabet and allowable lists do not establish an efficient decoder or a
polynomial-list improvement.

Recent limitations clarify this distinction. Alrabiah, Goyal and Guruswami
\cite[Theorems~1.1, 1.2 and~1.4]{alrabiahGoyalGuruswami2026} construct
super-polynomial lists just beyond Johnson for certain full-length
low-characteristic codes, and growing polynomial lists below capacity for
specified full-length prime-field codes. Their full length is the field
cardinality. Our family instead uses \(n\) subgroup points in
\(\mathbb F_{3^n}\), and its endpoint list \(\binom nk\) is allowed to
grow exponentially in \(n\) while meeting a fraction-of-field budget.
Neither that budget nor the separate pair-MCA event asserts polynomial-list
decoding up to capacity.

For exact finite thresholds, we use the Hamming grid and the worst-case
bad-challenge count
\[
  B_{\mathcal C}(r)=\max_p
  |\Bad_{\MCA}(\mathcal C,r/n,p)|.
\]
For the certified smooth prime-field instance
\(q=2^{128}+192513\equiv1\pmod {4096}\), \(n=4096\), and \(k=1024\), we prove
\(B_{\mathcal C}(0)=1\) and \(B_{\mathcal C}(1)\ge2\).  It follows that zero
is the greatest safe count radius and one the first unsafe count radius at
the \(2^{-128}\) MCA error target. Over real radii the safe interval is
\([0,1/4096)\), so its supremum is \(1/4096\) and is not itself safe.
The same code has greatest safe list count \(1536\) and first unsafe list
count \(1537\) at every positive interleaving width
(\cref{cor:boundary-field-list}).
On four length-\(64\) smooth codes with known exact MCA thresholds,
agreement packing gives the matching list transitions
(\cref{sec:jo-paired-thresholds}). The length-\(128\) subfield example
in \cref{cor:subfield-128-thresholds} has both transitions below capacity.
The prescribed-domain questions in \cite{arnonBonehFenzi2026} ask for
exact thresholds at rates \(1/2,1/4,1/8,1/16\). The general thresholds
remain open, including codes in the Johnson-to-capacity regime.

The coding theorems quantify over the complete ambient received-word and
affine-line spaces.  The simultaneous \(2^{-128}\) guarantee for capacity
lists and pair MCA at every radius uses the displayed exponential-size field.
An application with an independently defined source metric additionally
requires a comparison with the extracted pair-code distance.

\Cref{sec:polynomial-agreement} constructs deterministic certificates for
the whole-pair error count. Polynomial multipliers give a normalized linear
system whose first feasible response degree determines the strongest
certificate. For true count \(e\), the guarantee \(\min(e,n-k+1-e)\) is
exact through half the minimum distance, and explicit received pairs attain
the guarantee on each error shell of a nonzero domain.

The source application uses the cochain construction and observation profile
of \cite{vanDerVlugt2026protocol}. Its HCP and cellular Hodge primitives
are those of \cite{vanDerVlugt2026hcp}; the projection, reconstruction
and source-distance comparisons are proved here.
The coding theorems do not assume an HCP presentation or a commitment
protocol. A protocol using them must enforce the fixed-pair, later-challenge
order described in Section~6. Section~7 uses Paper 2's explicit systematic
cellular construction and proves its source-distance comparison. It retains
the original committed-cochain payload in a complete coordinate
layout and defines source validity intrinsically by an integral opening and
carry satisfying \(\delta_0a=z+qm\). Scalar cochain edits count \(2n\)
coordinates, while plane and RS edits count \(n\) whole pairs. The proved
global comparison is
\[
  d_C(z)\le d_P(z)=D_2(\bar z)\le2d_C(z).
\]
The factor two is necessary in the near direction; the far implication has
unit scale. An actual well-formed raw payload at
\((q,n,k)=(12289,4096,1024)\) has all three distances equal to \(3/4\).
This comparison uses modular source distances on raw payloads.
\Cref{prop:restricted-hodge-rs} supplies a separate full-Hodge comparison
by taking the native coupling matrix from the systematic RS graph.
With structured integer errors, reflected support and the unique-decoding
condition, a sitewise error bound controls the squared Hodge residual by the
pair-code distance. The large-field \(2^{-128}\) guarantee uses the field
specified above.
